
Items to follow up on from discussions:

\section{7/14}
\begin{itemize}
\item Lit review

  \begin{itemize}

  \item Generally: say more! Can spend e.g. a full page on how we differ from
    the flattening paper

  \item The lit review should stand alone (readers may start with it)

  \item Want to be superexplicit ``we pursue a more aggressive flattening
    strategy'' -- how? what makes it more agressive? what do we do that they
    don't?

  \item Avoid lumping stuff toether: e.g. the two ``serialized data papers,''
    the attempt to unify them as ``fast disk format'' is not really a match.
    Something like ``manual control of serialized formats'' is better

  \item \st{Don't use raw /cite -- the sentences should read like the footnote
    numbers don't exist. Use /citet to get better auto-formatat}

  \item To cite the Jane Street slide deck, best to try to make a bibtex entry
    that includes the talk as well

  \item \st{Consider using multiple /subsection or /paragraph blocks}

  \item Be sure to tie every cited paper exactly to what we did

  \item When contrasting, aim to emphasize fundamental differences of approach
    rather than just incidental implementation differences (e.g. we are using
    MLton and so we get some passes for free, but that doesn't make our strategy
    conceptually simpler)

  \end{itemize}

\item Tuple-flattening ``meat''

  \begin{itemize}

  \item Aim for text to be standalone! The english should stand alone. Be sure
    that we have concise, compact statements that describe what we do

  \item E.g. from ``the basic intution is'' -- just state ``the intuition is
    that we can optimize away matching pack/unpack operations by pushing them
    through function calls'' and then call out to diagrams that describe what we
    want to do

  \item Extend the example to show how the later \texttt{simplify} passes
    actually eliminate the tuple

  \item Consider a single, complex example rather than starting off with an easy
    case and adding complexity. We want it to be clear why we are necessary from
    the very start.

  \item For the algorithm part, we need to sketch it out more formally as
    transformations on an AST (in addition to perhaps an informal English
    description). Certainly the rewrite itself, maybe also the analysis (might
    be OK for that to be more informal)

  \item A useful way to write the AST transformation is ``[|foo bar = baz|] =
    [|foo bar|] = [|baz|], where ....'' followed by a detailed explanation. The
    transformation brackets can be subscripted by some symbol that indicates the
    result of the analysis 

  \end{itemize}

\end{itemize}

\section{6/29}
\begin{itemize}
\item Code snippets should be diagrams rather than inlined

\item A bulleted list should generally be compact on a single page

\item Introducing both C and python flavor pseudo-IRs up front is wordy and
  confusing. Also, using a C-flavor can be distracting in that it's unclear what
  we should ``interpret literally.'' Consider sticking to a single,
  Python-flavor IR

\item For the IR we do choose to pick, add an explicit grammar up-front

\item Add more ``signposting'' up-front to support hte nfo dump

\item \st{Probably worth adding a section on polyvariance / defunctionalization to
  motivate the con flattening}

\item With con flattening, maybe emphasize: MLton creates opportunites for
  specialization to happen downstream; we do the opposite -- explicitly
  specializing where we see a win

\item C struct definitions can be confusing because they imply stuff about
  alignment and packing that doesn't really hold

\item For diagrams: draw.io ?

\item Be careful to distinguish source language constructs from guarantees about
  generated code; e.g. ``tuples are allocated om the heap'' -- only when
  flattening fails

\item \st{Ok to introduce more sections -- e.g. separate for each optimization type}

\item Need to dig into why DeepFlatten is failing in our case or at least
  explicitly call out that we tried and couldn't do it. One useful thing to
  note: lattice-based ``poisoning'' analysis can be difficult to track

\item Add an intro where we make it more explicit what our contribution is

\item Think of the evaluation section as a reference to refer to throughout the
  other sections

\end{itemize}

Other items:
\begin{itemize}
\item We should avoid writing \texttt{'a} in SSA IR since it's monomorphic
\end{itemize}

%% \section{ 


